% Part: set-theory
% Chapter: spine
% Section: stagesbasics

\documentclass[../../../include/open-logic-section]{subfiles}

\begin{document}

\olfileid{sth}{spine}{Valphabasic}
\olsection{পর্যায়ের মৌলিক ধর্ম}

$V_\alpha$-গুলির সংজ্ঞা কেন ভিত্তিতত্ত্বের পক্ষে
গুরুত্বপূর্ণ, তা বোঝাতে কয়েকটি মৌলিক ফল দেখি।
প্রথমে একটি সংজ্ঞা:\footnote{``potent''-এর কোনো
প্রচলিত পরিভাষা নেই; \citet{ButtonLT1} এই নামটি
ব্যবহার করেছেন। এখানে ধর্মটি বোঝাতে
``উপসেট-অধোবদ্ধ'' বলছি।}
\begin{defn}
সেট $A$ \emph{উপসেট-অধোবদ্ধ} যদি এবং কেবল যদি
$\forall x((\exists y \in A)x \subseteq y \lif x \in A)$।
\end{defn}
\begin{lem}\ollabel{Valphabasicprops}
প্রত্যেক ক্রমসংখ্যা $\alpha$-এর জন্য:
\begin{enumerate}
	\item\ollabel{Valphatrans} প্রত্যেক $V_\alpha$ পরিযায়ী।
	\item\ollabel{Valphapotent} প্রত্যেক $V_\alpha$ উপসেট-অধোবদ্ধ।
	\item\ollabel{Valphacum} যদি $\gamma \in \alpha$ হয়, তবে $V_\gamma
	\in V_\alpha$ (আর তাই
	\olref{Valphatrans} অনুযায়ী
	$V_\gamma \subseteq V_\alpha$-ও)।
\end{enumerate}
\end{lem}

\begin{proof}
একসঙ্গে ট্রান্সফাইনাইট আবেশে ফলগুলি প্রমাণ
করি। আবেশের জন্য ধরি,
\olref{Valphatrans}--\olref{Valphacum} প্রত্যেক
$\beta < \alpha$ ক্রমসংখ্যার জন্য সত্য।

$\alpha = \emptyset$ ক্ষেত্রটি সরাসরি।

ধরি, $\alpha = \ordsucc{\beta}$।
\olref{Valphacum} দেখাতে $\gamma \in \alpha$
হলে আবেশ-অনুমান থেকে বা দুই সূচক সমান হওয়ায়
$V_\gamma \subseteq V_\beta$; তাই
$V_\gamma \in \Pow{V_\beta} = V_\alpha$।
\olref{Valphapotent} দেখাতে ধরি,
$A \subseteq B \in V_\alpha$; অর্থাৎ
$A \subseteq B \subseteq V_\beta$। ফলে
$A \subseteq V_\beta$, তাই $A \in
V_\alpha$। \olref{Valphatrans} দেখাতে
লক্ষ করি, $x \in A \in V_\alpha$ হলে
$A \subseteq V_\beta$; তাই $x \in V_\beta$।
আবেশ-অনুমানে $V_\beta$ পরিযায়ী বলে
$x \subseteq V_\beta$; অতএব $x \in V_\alpha$।

এবার ধরি, $\alpha$ একটি সীমা ক্রমসংখ্যা।
\olref{Valphacum} দেখাতে $\gamma \in \alpha$
হলে $\gamma \in \ordsucc{\gamma} \in \alpha$।
আবেশ-অনুমানে $V_\gamma \in
V_{\ordsucc{\gamma}}$; তাই
$V_\gamma \in \bigcup_{\beta \in \alpha} V_\beta = V_\alpha$।
\olref{Valphatrans} ও \olref{Valphapotent}
দেখাতে শুধু লক্ষ করলেই চলে: পরিযায়ী
সেটগুলির সংযোগ পরিযায়ী, আর
উপসেট-অধোবদ্ধ সেটগুলির সংযোগও
উপসেট-অধোবদ্ধ।
\end{proof}

\begin{lem}\ollabel{Valphanotref}
প্রত্যেক ক্রমসংখ্যা $\alpha$-এর জন্য
$V_\alpha \notin V_\alpha$।
\end{lem}

\begin{proof}
ট্রান্সফাইনাইট আবেশে প্রমাণ করি। স্পষ্টতই
$V_\emptyset \notin V_\emptyset$।

যদি $V_{\ordsucc{\alpha}} \in V_{\ordsucc{\alpha}} = \Pow{V_\alpha}$
হয়, তবে $V_{\ordsucc{\alpha}} \subseteq V_\alpha$।
আবার \olref{Valphabasicprops} অনুযায়ী
$V_\alpha \in V_{\ordsucc{\alpha}}$; তাই
$V_\alpha \in V_\alpha$। প্রতিবিপরীত রূপে,
$V_\alpha \notin V_\alpha$ হলে
$V_{\ordsucc{\alpha}} \notin V_{\ordsucc{\alpha}}$।

$\alpha$ সীমা হলে এবং $V_\alpha \in V_\alpha = \bigcup_{\beta \in
\alpha}V_\beta$ হলে কোনো $\beta \in
\alpha$-এর জন্য $V_\alpha \in V_\beta$।
কিন্তু \olref{Valphabasicprops} দুবার
প্রয়োগে $V_\beta \in V_\alpha$-ও; ফলে
$V_\beta \in V_\beta$। প্রতিবিপরীত রূপে,
প্রত্যেক $\beta \in \alpha$-এর জন্য
$V_\beta \notin V_\beta$ হলে
$V_\alpha \notin
V_\alpha$।
\end{proof}

\begin{cor}
যেকোনো ক্রমসংখ্যা $\alpha, \beta$-র জন্য:
$\alpha \in \beta$ যদি এবং কেবল যদি
$V_\alpha \in V_\beta$।
\end{cor}

\begin{proof}
এক দিক \Olref{Valphabasicprops} থেকে পাই।
বিপরীত দিকে ধরি, $V_\alpha \in V_\beta$।
\olref{Valphanotref} অনুযায়ী
$\alpha \neq \beta$। আর $\beta \notin \alpha$;
নতুবা $V_\beta \in V_\alpha$ হতো এবং
\olref{Valphabasicprops} দুবার প্রয়োগে
$V_\beta \in V_\beta$ পেতাম,
যা \olref{Valphanotref}-এর বিরোধী।
তাই ত্রিবিভাজন থেকে $\alpha \in \beta$।
\end{proof}
\noindent
এগুলি থেকে প্রত্যেক $V_\alpha$-কে স্তরক্রমের
$\alpha$-তম পর্যায় হিসেবে ভাবা যায়। কারণটি দেখি।

$V_\alpha$-গুলি যে একটি \emph{পর্যায়ক্রমিক}
প্রক্রিয়ায় গঠিত, তা ভাবা যায়: ক্রমসংখ্যার
ব্যবহার পুনরাবৃত্তির ধারণার সঙ্গে মেলে। আবার
এক পর্যায় অন্যটির আগে গঠিত হলে, অর্থাৎ
$V_\beta \in V_\alpha$, অর্থাৎ $\beta \in \alpha$
হলে, প্রক্রিয়াটি \emph{সঞ্চয়ী}; কারণ
$V_\beta \subseteq V_\alpha$। সবশেষে, আগের
যেকোনো পর্যায়ে পাওয়া সেটগুলির \emph{সমস্ত}
সম্ভব সমষ্টিও আমরা গঠন করছি: উত্তরসূরি পর্যায়
$V_{\ordsucc{\alpha}}$ হলো পূর্ববর্তী
$V_\alpha$-র ঘাত সেট।

সারকথা, $\ZFminus$ দিয়ে সেটের সঞ্চয়ী-পর্যায়ক্রমিক
স্তরক্রমের ধারণা প্রকাশের কাজ \emph{প্রায়}
শেষ। তবে প্রতিস্থাপনের ন্যায্যতা এখনও দিতে হবে।

\end{document}
